%LTeX: language=fr-FR
\documentclass[12pt,a4paper]{article}
\usepackage[utf8]{inputenc}
\usepackage[T1]{fontenc}
\usepackage[french]{babel}
\usepackage[margin=1.5cm]{geometry}
\usepackage{amsmath,amssymb}
\usepackage{graphicx}
\usepackage{enumitem}
\usepackage{tcolorbox}
\usepackage{xcolor}
\usepackage{multicol}
\usepackage{tasks}
\usepackage{tikz}
\usetikzlibrary{arrows.meta}

\tcbuselibrary{skins}

\definecolor{bandeau}{HTML}{E8EEF7}
\definecolor{titrebleu}{HTML}{1A4FA0}

% Bandeau de titre pour chaque exercice
\newtcbox{\exobox}{
  colback=bandeau, colframe=bandeau, boxrule=0pt,
  arc=0pt, left=4pt, right=4pt, top=3pt, bottom=3pt,
  fontupper=\bfseries, width=\linewidth, halign=left
}

% \exercice{numéro}{titre}{barème}
\newcommand{\exercice}[3]{%
  \noindent\exobox{Exercice #1 --- #2 \hfill (#3 points)}\par
}

\pagestyle{empty}
\setlist{topsep=2pt, itemsep=3pt, parsep=0pt}
\setlength{\parskip}{2pt}
\setlength{\parindent}{0pt}

% Droite graduée de -10 à 10 (même droite que l'exercice 6 de la feuille d'exercices)
\newcommand{\droitecinq}{%
  \begin{tikzpicture}[x=0.7cm,baseline=-0.5ex]
    \draw[{Latex[length=2mm]}-{Latex[length=2mm]}] (-10.6,0) -- (10.6,0);
    \foreach \i in {-10,...,10} {
      \draw (\i,0.08) -- (\i,-0.08);
      \node[below] at (\i,-0.08) {\scriptsize$\i$};
    }
  \end{tikzpicture}%
}

% Un cas de l'exercice 4 : énoncé + droite + zone de réponse
\newcommand{\casinter}[2]{%
  \par\noindent
  \begin{minipage}{\linewidth}
    #1. $#2$\par
    \begin{center}
      \droitecinq\\[0.5em]
      $I\cap J=$ \makebox[3.5cm]{\dotfill} \qquad\qquad $I\cup J=$ \makebox[3.5cm]{\dotfill}
    \end{center}
  \end{minipage}\par
}

\begin{document}

\begin{center}
  {\Large\bfseries Contrôle --- Ensembles de nombres (1)~: intervalles}\\[0.4em]
  {\itshape Seconde \quad $\bullet$ \quad Durée~: 50 minutes \quad $\bullet$ \quad Calculatrice interdite}
\end{center}

\vspace{0.4em}
\noindent Nom et prénom~: \dotfill

\vspace{0.3em}
\noindent\textit{Exercices 1 et 4 \underline{sur la copie}, exercices 2 et 3 \underline{sur le sujet} (à rendre avec la copie). La rédaction et la justification des réponses sont prises en compte dans la notation.}

% ================= Exercice 1 =================
\exercice{1}{Automatismes}{5}
\underline{Sur la copie.} Les questions sont indépendantes. Justifier chaque réponse.
\begin{enumerate}[itemsep=6pt]
  \item Calculer $2^{5}\times 2^{-3}$.
  \item On sait que $\dfrac{32}{33}-\dfrac{31}{32}>0$. Compléter avec le symbole $<$ ou $>$ : \quad $\dfrac{31}{32}\ \dots\ \dfrac{32}{33}$.
  \item On considère l'expression $T=M(5+a)$. Calculer la valeur de $T$ pour $M=4$ et $a=\dfrac{3}{4}$.
  \item Factoriser à l'aide d'une identité remarquable : \quad $A=x^2-6x+9$
\end{enumerate}

\vfill
% ================= Exercice 2 =================
\exercice{2}{Intervalles}{5}
\underline{Sur le sujet.}
\begin{enumerate}
  \item Compléter les pointillés par $\in$ ou $\notin$.

  \quad a) $4\ \dots\ [-1\,;\,4[$ \hfill b) $-3\ \dots\ ]-5\,;\,-3]$ \hfill c) $\dfrac{2}{3}\ \dots\ [0{,}5\,;\,1]$ \hspace*{2cm}
  \item Compléter le tableau.
  \begin{center}
  \renewcommand{\arraystretch}{1.8}
  \begin{tabular}{|c|c|c|c|c|}
  \hline
  Inégalité & $-3\leqslant x<4$ & $x>7$ & \makebox[3cm]{\dotfill} & \makebox[3cm]{\dotfill} \\ \hline
  Intervalle & $x\in\dots\dots\dots\dots$ & $x\in\dots\dots\dots\dots$ & $x\in\,]-2\,;\,5]$ & $x\in\,]-\infty\,;\,2{,}5[$ \\ \hline
  \end{tabular}
  \end{center}
\end{enumerate}

\vfill
% ================= Exercice 3 =================
\exercice{3}{Intersection et union}{4}
\underline{Sur le sujet.} Représenter $I$ et $J$ sur la droite en deux couleurs, puis donner $I\cap J$ et $I\cup J$.

\casinter{1}{I=\,]-\infty\,;\,1]\ ;\ J=[-4\,;\,4[}
\casinter{2}{I=[-5\,;\,-1[\ ;\ J=[2\,;\,+\infty[}

\vfill
% ================= Exercice 4 =================
\exercice{4}{Inégalités et inéquations}{6}
\underline{Sur la copie.}
\begin{enumerate}
  \item Résoudre les inéquations. Donner les solutions sous forme d'intervalle.

  \quad a) $-4x+1>13$ \qquad\qquad b) $7-2x<x+1$
  \item Compléter en justifiant par une suite d'inégalités : \quad $x\in\,]-1\,;\,3] \iff -2x+4\in\ \dots\dots$
\end{enumerate}

\vfill
\end{document}
